%!TEX root = ../hutbthesis_main.tex

\chapter{系统测试与精度分析}
\section{测试目的与方案}
本章在 User-in-the-Box 官方框架的基础上，对整个系统进行综合测试，测试会用到框架自带的方向盘模型、上肢肌肉骨骼模型和控制接口，用来检验肌肉驱动生物力学人机交互系统的功能是否完整、控制是否精准、运行是否稳定以及运动是否符合人体规律，测试会按照模型加载、肌肉驱动、轨迹跟踪、数据采集、精度评估的固定步骤来进行，测试使用统一的实验方法和量化指标，保证得到的结果真实、客观且可以重复实现，所有测试工作都在本地的项目环境中完成，也能为课题研究提供可信的实验依据。

本次测试主要是为了检验系统能不能达到毕业设计要求的功能和性能，重点检查肌肉驱动控制方向盘的精准程度和稳定程度，同时也检验系统各个模块一起工作是否可靠，测试围绕功能完整性、控制精度、运行稳定性、运动生物合理性、任务完成效果这五个主要目标开展，借助 uitb 框架自带的方向盘模型和 SteeringTask 任务接口，搭建出肌肉驱动、方向盘转动、精度评估的完整测试流程。

测试使用固定参数、多种情况对比的方式，统一仿真环境和控制参数，减少因为模型不同、参数变化带来的结果误差，测试内容包含数据处理、运动学计算、肌肉驱动控制、视觉引导、仿真运行、指标评估等全部模块，可以全面体现系统的实际性能，测试得到的结果会直接用在论文的数据说明和结论证明中。
\section{测试环境与参数}
本次测试基于 uitb 官方框架搭建的仿真环境，软硬件配置如下，确保仿真运行流畅、数据采集准确：

硬件环境：常规PC平台（满足 MuJoCo 仿真运行需求，无特殊硬件要求）；

软件环境：Windows 10 64位操作系统、Python 3.8、MuJoCo 2.3、User-in-the-Box 框架（原生未修改）；

开发工具：VS Code、NumPy、Pandas、Matplotlib（用于数据处理与结果可视化）；

项目结构：依托 User-in-the-Box 框架原生目录结构，调用内置方向盘模型与任务接口。

为保证测试的一致性与可复现性，统一设置以下核心参数，与论文设计参数完全一致：

仿真步长：0.002s（与 uitb 框架默认仿真步长一致）；

方向盘转动范围：$±90°$（调用框架原生方向盘关节约束，关节名：steering\_joint）；

测试时长：单工况测试60s，稳定性测试3600s（1小时）；

控制策略：PPO强化学习策略+自适应映射算法（与第五章算法设计一致）；

目标轨迹：阶跃轨迹（定点转向）、正弦轨迹（连续转向），模拟实际方向盘操控场景。
\section{测试模块与指标体系}
系统测试严格按照五大模块、四类评估维度开展，具体指标如下表所示。
\begin{table}[htbp]
  \centering
  \caption{系统测试模块与评估指标}
  \label{tab:test_index}
  \begin{tabular}{p{2.5cm}p{5cm}p{7cm}}
    \toprule
    测试模块 & 主要测试内容 & 核心评价指标 \\
    \midrule
    数据处理模块 & 数据加载、滤波、标准化、时序对齐 & 数据同步误差、预处理延迟 \\
    运动学模块 & SMPL 前向运动学解算、3D 姿态输出 & 关节位置误差、运动平滑度、关节超限率 \\
    控制模块 & 角度跟踪、响应速度、控制稳定性 & 平均绝对误差 MAE、均方根误差 RMSE、响应延迟、最大误差 \\
    仿真模块 & 肌肉骨骼动力学、方向盘物理仿真 & 仿真帧率、无故障运行时间、角度漂移量 \\
    评估模块 & 轨迹相似度、生物合理性、任务完成度 & 跟踪成功率、任务完成率、整体稳定性评分 \\
    \bottomrule
  \end{tabular}
\end{table}

\section{功能测试}
为了检验肌肉和方向盘耦合驱动系统的功能是否完整、运行是否稳定以及操作是否好用，本章将依托已经搭建好的仿真平台和肌肉驱动的上肢模型，设计三类常见的人机交互任务来开展系统功能测试。

这三类任务分别是遥控驾驶任务、控制精度任务和小物体触碰任务，通过这三类不同的任务，全面检查模型在不同交互场景下的运动反应、控制精准度以及动态稳定效果。

测试工作需要借助 User-in-the-Box 框架自带的统一启动脚本 hmi\_quickstart.py 来开展，通过这个脚本能够实现不同测试任务的快速切换。就像图~\ref{fig:test_scheme}的测试方案与指标体系图所展示的那样，工作人员可以通过命令行参数，自主选择人工操控模式或者自动评估模式。

其中，人工操控模式支持实时控制操作，自动评估模式则可以实时记录相关数据、计算评估指标，同时还能将测试数据导出保存，确保整个测试过程规范统一，测试结果能够用具体数据来体现，避免主观判断。

这种测试方式既保证了测试流程的标准化，也让测试结果更加客观可量化，能够全面检验肌肉–方向盘耦合驱动系统的实际使用效果，这样系统在不同场景下能够稳定运行、满足使用需求。
\begin{figure}[htbp]
    \centering
    \includegraphics[width=0.8\textwidth]{image_11.png}
    \caption{测试方案与指标体系}
    \label{fig:test_scheme}     
\end{figure}

功能测试用于验证系统能否正常启动、加载、运行、控制与评估，是系统可用的基础。下图~\ref{fig:function_test}是User-in-the-Box 框架仿真环境搭建成功的画面，正在进行测试
\begin{figure}[htbp]
    \centering
    \includegraphics[width=0.8\textwidth]{image_12.png}
    \caption{环境搭建成功画面}
    \label{fig:function_test}       
\end{figure}

在开展系统测试时，会分三类任务进行测试，每类任务对应不同的脚本、场景和测试重点，具体如下：在遥控驾驶任务中，系统会调用脚本文件来进行控制，同时加载 mobl\_arms\_index\_remote\_driving 仿真场景，采用 WASD 按键操作来模拟摇杆推拉指令，进而实现对上肢模型和方向盘的连续转向控制，这项测试主要是为了检验系统在长时间连续操作下，是否能保持响应流畅、姿态稳定，不会出现卡顿问题。测试期间，会保持肌肉驱动相关参数不变，同步记录方向盘转动角度、上肢关节的运动轨迹，以及仿真实时运行情况，以此验证系统在模拟驾驶类连续操控场景中的实际使用效果。

在控制精度测试任务中，系统会加载 \texttt{pointing\_teleop.py} 控制脚本，同时启用$ \texttt{mobl\_arms\_index\_pointing} $高精度场景。工作人员通过 WASD 按键发出控制指令，驱动上肢模型的手指末端对准固定目标点。重点检验肌肉驱动模型的定位准确性、细微动作的控制能力，以及姿态调整的速度。通过记录目标偏差距离、稳定运行时间等具体数据，来判断系统在精细操作场景下的控制效果。

在小物体触碰任务中，会启用控制脚本文件来和场景文件来进行控制和紧凑场景，工作人员通过按键发出控制指令，驱动上肢模型的手指完成对小目标物体的精准触碰操作，重点测试模型在狭小空间内的运动灵活性、躲避障碍物的合理性，以及触碰时的稳定性，以此验证肌肉与方向盘耦合驱动系统在复杂约束环境下，是否能稳定运行、满足使用需求。

此外，下图~\ref{fig:accuracy_test}中，还展示了肌肉驱动方式下，方向盘角度跟踪精度的对比数据，可直观体现系统的控制精度情况。
\begin{figure}[htbp]
    \centering
    \includegraphics[width=0.8\textwidth]{image_13.png}
    \caption{肌肉驱动方向盘角度跟踪精度对比}
    \label{fig:accuracy_test}   
\end{figure}    

三种测试任务都可以使用人工控制和自动评价两种运行方式，在程序中加入 --eval 参数，就能开启自动检测流程，系统会载入模型相关参数，按照设定的次数完成自动测试，并输出可以量化的评价数据，加入 --csv 参数，就能在人工操作时实时导出动作数据、控制偏差和运行稳定程度，这些数据可以让测试结果更加真实可信。测试结果可以说明，肌肉与方向盘结合的驱动系统，在三种任务里都可以保持稳定的运行状态，上肢模型动作流畅不会出现晃动，方向盘可以和上肢同步动作，不会出现时间滞后，关节活动范围合理，不会超出限制，位置确定精度和物体触碰成功率，都可以达到人机交互的使用标准。遥控驾驶任务中，系统可以持续运行不会出现故障，转向反应速度快，模型姿态保持稳定，控制精度任务中，目标位置的偏差数值较小，能快速达到稳定状态，小物体触碰任务中，模型动作灵活，触碰位置准确，可以在空间条件有限的情况下完成指定操作。综合所有测试结果可以看出，本次搭建的肌肉驱动上肢模型和结合驱动系统，设计结构合理，运行状态稳定，功能模块完善，可以满足模拟驾驶、精细操作、交互触碰等多种场景的使用需要，具备较好的实际使用价值和拓展空间，图~\ref{fig:tracking_curve}为角度跟踪曲线数据，图~\ref{fig:step_tracking_curve}为阶跃轨迹下方向盘角度跟踪曲线数据。
\begin{figure}[H]
    \centering
    \includegraphics[width=0.8\textwidth]{image_14.png}
    \caption{角度跟踪曲线数据}
    \label{fig:tracking_curve}
\end{figure}

\begin{figure}[H]
    \centering
    \includegraphics[width=0.8\textwidth]{image_15.png}
    \caption{阶跃轨迹下方向盘角度跟踪曲线数据}
    \label{fig:step_tracking_curve}     
\end{figure}













为验证肌肉驱动人机交互系统各功能模块的正确性与可靠性，开展单元测试。测试覆盖模型加载、关节驱动、力矩映射、信号滤波、死区抑制与限幅保护六大核心模块，结果如图~\ref{fig:unit_test}所示：
\begin{figure}[H]
    \centering
    \includegraphics[width=0.8\textwidth]{image_16.png}
    \caption{单元效果测试图}
    \label{fig:unit_test}   
\end{figure}











通过对系统的多个部分进行单独测试，以此确认每个模块都可以正常使用，模型加载测试会对肌肉骨骼模型的初始运行状态进行检查，模型在启动过程中不会出现错误提示，也不会发生异常情况，模型可以被顺利加载到系统当中，关节驱动测试会使用连续变化的运动方式开展检验，关节可以按照设定的轨迹平稳进行活动，关节角度的变化过程不会出现明显的异常，整体运行状态保持稳定，力矩映射测试会对肌肉信号与输出动力之间的关系进行验证，数据的分布情况和预期的变化规律基本保持一致，信号到动力的转换过程不会出现偏差，信号滤波测试会对处理前和处理后的信号进行对比，经过滤波处理的信号可以减少杂乱波动的影响，信号的整体质量会得到明显改善，死区抑制测试会对信号在零值附近的波动情况进行检验，波动问题可以得到有效控制，系统在静止状态下不会出现错误动作，限幅保护测试会对输出信号的范围进行约束，超出安全范围的信号会被限制在合理区间内，系统和相关设备可以得到有效保护。通过以上各项测试可以发现，系统中的每个功能模块都可以正常运行，模块的运行效果可以达到设计时的相关要求，这些单独测试的顺利完成，也为后续的整体组合测试和系统综合测试打下了扎实的运行基础。
